% Generated by analysis_for_adversarial_report.Rmd; do not edit numerical cells.
\begin{table}[!htbp]
\singlespacing\renewcommand{\arraystretch}{1.2}
\centering
\caption{Entering-Class Composition at Ninth-Grade Entry}\label{tab: ninth grade composition}
\small
\begin{tabular*}{\textwidth}{@{\extracolsep{\fill}}p{0.46\textwidth}cc@{}}
\toprule
Specification & \shortstack{Prior-year GPA\\(SD)} & \shortstack{Low-income share\\(pp)} \\
 & (1) & (2) \\
\midrule
All grades & -0.149 & 3.50 \\
 & (0.032) & (0.92) \\
95\% confidence interval & [-0.212, -0.086] & [1.71, 5.30] \\
School-grade-years & 57,269 & 72,490 \\
\midrule
Ninth grade & -0.136 & 4.41 \\
 & (0.063) & (2.34) \\
95\% confidence interval & [-0.259, -0.013] & [-0.19, 9.00] \\
School-grade-years & 7,594 & 7,607 \\
\bottomrule
\end{tabular*}
\par\smallskip
\begin{minipage}{\textwidth}\footnotesize
Notes: The outcomes and estimator follow Table~\ref{tab: alloc headline}. The ninth-grade specification restricts the sample to ninth-grade entrants. Both specifications use not-yet-treated school-grades, a universal pre-adoption base period, no additional covariates, and the unbalanced-panel estimator described in Appendix~\ref{app:access}. Within comparisons, school-grade cells receive equal weight; group-time effects are aggregated by adoption-group size. The all-grade and ninth-grade estimates cover supported post-adoption entering classes in 2017--2020 and 2017--2019, respectively. For ninth grade, the Metropolitan Region supplies comparisons but has no identified post-adoption effect. Observation counts include pre-adoption years from 2012. Region-clustered SE from 10,000 multiplier-bootstrap draws are in parentheses; brackets report pointwise 95\% intervals. All-grade point estimates reproduce the main table; SE are recomputed with fresh draws. No permutation-test significance stars are reported. The small number of regions remains a limitation for inference.
\end{minipage}
\end{table}
